\subsection{13.42: a tensor completion that is not locally nilpotent}
\label{cand:13.42}
Let \(A=\Z[t]\) and \(N=F(x,y)/\gamma_4F(x,y)\).
Its unrestricted tensor \(A\)-completion is not locally nilpotent.
The Notebook refers to Myasnikov--Remeslennikov's foundational paper
\cite{exponential-groups-1}. Here we use the power, conjugation,
commuting-product axioms and tensor-completion universal property
as restated in \cite[Definitions~1--2]{exponential-groups}.
The universal category allows scalar torsion and does not restrict
targets to nilpotent groups. This distinction is part of the claim.

Let \(L\) be the rational BCH group of the Lie algebra with basis
\(X,Y,U,V,W\), brackets
\[
 [X,Y]=U,\qquad [X,U]=V,\qquad [Y,U]=W,
\]
and all other basic brackets zero.
Coordinates \((a,b,c,d,e)\) denote
\(\exp(aX+bY+cU+dV+eW)\), and \(\beta(a,b,c,d,e)=b\) is a
homomorphism \(L\to(\Q,+)\).
These are finite polynomial operations: realize the Lie algebra in
the free associative algebra on \(X,Y\) truncated above degree three
and use \(\exp Z=1+Z+Z^2/2+Z^3/6\).
For \(l=(a,b,c,d,e)\), \(h=(A,B,C,D,E)\), BCH gives the first
three coordinates of \(l^h=h^{-1}lh\) as
\[
 (a,b,c+aB-bA).
\]
It also gives the commuting criterion
\[
 [l,h]=1\quad\Longleftrightarrow\quad
 aB-bA=aC-cA=bC-cB=0.
\]
For example, equality of the degree-two coordinates in the two
products first gives \(aB-bA=0\); then equality in degree three
gives the other two minors. Thus commuting elements have proportional
triples \((a,b,c)\), \((A,B,C)\), including the zero cases.

Let \(M=\Q[\Q]\) additively, with basis \(T^q\), \(q\in\Q\).
Give it the right \(L\)-action \(v^l=vT^{\beta(l)}\), and define
\[
 F(a,b,c,d,e)=
 \begin{cases}aT^{c/a}&b=0,\ a\ne0,\\0&\text{otherwise}.\end{cases}
\]
The conjugation formula, coordinate powers, and proportional-triple
criterion prove, respectively,
\[
 \begin{aligned}
 F(l^h)&=F(l)T^{\beta(h)},&
 F(l^n)&=nF(l),\\
 F(lh)&=F(l)+F(h)\quad([l,h]=1),&
 F(l)T^{\beta(l)}&=F(l).
 \end{aligned}
\]
For the third assertion, the defining function of a triple is
homogeneous over \(\Q\), hence additive on any line.
It is not asserted to be additive on arbitrary pairs.

Form the right-action semidirect product with multiplication
\[
 H=L\ltimes M,\qquad (l,v)(h,w)=(lh,vT^{\beta(h)}+w),
\]
write \(i(v)=(1,v)\), and put \(f(l,v)=i(F(l))\).
Then
\[
 f(g^h)=f(g)^h,\quad [g,f(g)]=1,\quad
 f(gh)=f(g)f(h)\ \text{if }[g,h]=1,
\]
and, for all integers \(n,m\),
\[
 f(g^nf(g)^m)=f(g)^n.
\]
The last identity follows by projecting to \(L\).
If \(g,h\) commute, equivariance also gives
\([f(g),h]=[g,f(h)]=1\).

For \(P(t)\in\Z[t]\), define
\[
 g^P=g^{P(0)}f(g)^{P'(0)}.
\]
This supplies all the \(A\)-group axioms. Unit, zero and addition
follow from \([g,f(g)]=1\).
If \(P(0)=a,P'(0)=b,Q(0)=c,Q'(0)=d\), then
\[
 (g^P)^Q=g^{ac}f(g)^{bc+ad}=g^{PQ},
\]
using \(f(g^P)=f(g)^a\).
Equivariance proves the conjugation axiom. For commuting \(g,h\),
the cross-commutations just established allow rearrangement to give
\((gh)^P=g^Ph^P\).
Thus \(H\) is an \(A\)-group in the required unrestricted category.
It is torsion-free as an ordinary group, being an extension of
torsion-free groups, although its scalar action factors through
\(A/(t^2)\).

Set \(x=(\exp X,0)\), \(y=(\exp Y,0)\).
They generate a group of class at most three, so the map from the
free nilpotent group \(N\) extends by the tensor-completion universal
property to an \(A\)-homomorphism \(N^A\to H\).
But \(x^t=i(T^0)\), and direct multiplication gives
\[
 [i(v),y]=i(v(T-1)).
\]
For every \(r\ge1\), the image of \([x^t,y,\ldots,y]\), with \(r\)
copies of \(y\), is \(i((T-1)^r)\ne1\): its coefficient at \(T^r\)
is one. The fixed pair \(x^t,y\) therefore generates a nonnilpotent
subgroup of \(N^A\), proving the claim.
The map need not be onto \(H\), and injectivity of \(N\to H\) is
unnecessary. Completion within nilpotent \(A\)-groups would be a
different construction. See \artifact{research/13.42-proof.md}.
